
\subsubsection{FFGraph}

The class {\bf FFGraph} is a container for the class {\bf
  FFGraphClass}. As the user has no direct access to {\bf
  FFGraphClass}, using {\bf FFGraph} is the only way to call
methods in {\bf FFGraphClass} objects. This is done by using the
\verb|->| operator. E.g. to add a node to a graph you would write:

\bigskip

{\tt
FFGraph g;\\
g\verb|->|addNode();\\
}

\bigskip

FFGraph does an automatic garbage collection. This means that if
there is no container for a graph left, the graph is freed.

FFGraph defines the following methods:

\begin{itemize}

\item{\bf FFGraph()}
  
  Creates a new graph. The underlying {\bf FFGraphClass} object
  is not allocated before it is actually used.

\item{\bf FFGraph\& operator = (const FFGraph\& graph)}
  
  Sets the object to be a container for the same graph as {\bf
    graph}.

\item{\bf int operator == (const FFGraph\& graph)}
  
  Returns whether the object is a container for the same graph as
  {\bf graph}.

\item{\bf int operator != (const FFGraph\& graph)}
  
  Returns whether the object is not a container for the same
  graph as {\bf graph}.
  
\item{\bf FFGraphClass * operator \verb|->| () const}
  
  Passes the method to the underlying {\bf FFGraphClass} object.

\end{itemize}

\subsubsection{FFGraphClass}

The class {\bf FFGraphClass} provides standard methods to
manipulate graphs. The methods of {\bf FFGraphClass} can not be
accessed directly, but must be accessed using the class {\bf
  FFGraph}.

{\bf FFGraphClass} defines the following methods:

\begin{itemize}

\item{\bf BOOL getDirected()}
  
  Returns whether the graph is directed. This method is intended
  to be used in the future, as only directed graphs are supported
  at the moment. May return any value.

\item{\bf int getGraphicNodeIndex()}
  
  If the nodes of the graph are labeled with {\bf FFGraphicNode},
  the index of the label is returned. If not, -1 is returned.

\item{\bf int getGraphicEdgeIndex()}
  
  If the edges of the graph are labeled with {\bf FFGraphicEdge},
  the index of the label is returned. If not, -1 is returned.

\item{\bf int getGraphWindowIndex()}
  
  If the graph is labeled with a window label, the index of the
  label is returned. If not, -1 is returned.
 
% inline void setNodeIndex(int i) {nodeIndex=i;}
% inline void setEdgeIndex(int i) {edgeIndex=i;}
% inline void setGraphIndex(int i) {graphIndex=i;}
  
\item{\bf void load(const String\& buffer, BOOL datei)}
  
  If {\bf datei == YES} the graph specified in the sgraph
  compatible file {\bf buffer}
  is merged into the graph.\\
  If {\bf datei == NO} the graph that is contained in {\bf
    buffer} in sgraph format is merged into the graph.

\item{\bf void saveToFile(const String\& file)}
  
  The graph is saved to the file {\bf file}.
  
\item{\bf void saveToFile(const String\& file, const
    AVLSet\verb|<|FFNodeClass *\verb|>|\& set)}
  
  The subgraph defined by the nodes in {\bf set} and the edges
  between the nodes in {\bf edge} is saved to file {\bf file}.

\item{\bf void saveToString( String\& str)}

The graph is saved to the string {\bf str}.

\item{\bf void saveToString( String\& str, const
    AVLSet\verb|<|FFNodeClass *\verb|>|\& set)}
      
  The subgraph defined by the nodes in {\bf set} and the edges
  between the nodes in {\bf edge} is saved to string {\bf str}.

\item{\bf FFNode addNode(const FFPoint\& p=FFPoint(100,100,0))}
  
  A new node is added to the graph and returned. If the nodes of
  the graph are labeled with {\bf FFGraphicNode}, the node is
  moved to position {\bf p}.

\item{\bf  FFEdge addEdge(FFNode s, FFNode t)}
  
  A new edge between source node {\bf s} and target node {\bf t}
  is added to the graph and returned. {\bf s} and {\bf t} must
  belong to this graph.

\item{\bf void removeNode(FFNode n)}
  
  Node {\bf n} is removed from the graph. All edges from or to
  {\bf n} are removed from the graph as well.

\item{\bf void removeEdge(FFEdge e)}

Edge {\bf e} is removed from the graph.

\item{\bf void clear()}

All nodes and edges are removed from the graph.

\item{\bf void allNodes(AVLSet\verb|<|FFNodeClass *\verb|>|\& set)}

All nodes of the graph are added to set {\bf set}.

\item{\bf void allEdges(AVLSet\verb|<|FFEdgeClass *\verb|>|\& set)}

All edges of the graph are added to set {\bf set}.

\item{\bf const AVLSet\verb|<|FFNodeClass *\verb|>|\& allNodes()
    const}

A set containing all nodes of the graph is returned.

\item{\bf const AVLSet\verb|<|FFEdgeClass *\verb|>|\& allEdges()
    const}

A set containing all edges of the graph is returned.

\item{\bf int getNodeLabelIndex(const String\& type,
    AVLSet\verb|<|int\verb|>| *set=NULL) const}
  
  If {\bf set == NULL} the index of the node label with type {\bf
    type} is returned, or -1 if none is found.  if {\bf set !=
    NULL} the indexes of all node labels which contain {\bf type}
  in their type are added to {\bf set}. The index of the first
  label that is found is returned, or -1, if none is found.
  
\item{\bf int getEdgeLabelIndex(const String\&,
    AVLSet\verb|<|int\verb|>| *set=NULL) const}
  
  If {\bf set == NULL} the index of the edge label with type {\bf
    type} is returned, or -1 if none is found.  if {\bf set !=
    NULL} the indexes of all edge labels which contain {\bf type}
  in their type are added to {\bf set}. The index of the first
  label that is found is returned, or -1, if none is found.
  
\item{\bf int getGraphLabelIndex(const String\&,
    AVLSet\verb|<|int\verb|>|* set=NULL) const}
  
  If {\bf set == NULL} the index of the graph label with type
  {\bf type} is returned, or -1 if none is found.  if {\bf set !=
    NULL} the indexes of all graph labels which contain {\bf
    type} in their type are added to {\bf set}. The index of the
  first label that is found is returned, or -1, if none is found.
  
\item{\bf FFGraphLabel *getGraphLabel(const String\& s,
    AVLSet\verb|<|FFGraphLabel *\verb|>| *set=NULL)}
  
  If {\bf set == NULL} the graph label with type {\bf type}
  is returned, or {\bf NULL} if none is found.  if {\bf set !=
    NULL} all graph labels which contain {\bf type} in their type
  are added to {\bf set}. The first label that is found is
  returned, or {\bf NULL}, if none is found.

\item{\bf FFGraphLabel *getGraphLabel(int index)}
  
  The graph label with index {\bf index} is returned or {\bf
    NULL} if the graph has no graph label with index {\bf index}.
  
\item{\bf const MPlex\verb|<|FFGraphLabel *\verb|>|\&
    getGraphLabels() const}

An MPlex containing all graph labels is returned.

\item{\bf const MPlex\verb|<|FFNodeLabel *\verb|>|\&
    getNodeLabels() const}

An MPlex containing the prototypes for all node labels is returned.

\item{\bf const MPlex\verb|<|FFEdgeLabel *\verb|>|\&
    getEdgeLabels() const}

An MPlex containing the prototypes for all edge labels is returned.

\item{\bf int addNodeLabel(FFNodeLabel *label)}
  
  Prototypes of label {\bf label} are added to all nodes in the
  graph. The index of the label is returned. If the nodes are
  already labeled with a label {\bf label}, only the index is
  returned.

\item{\bf int addEdgeLabel(FFEdgeLabel *label)}
  
  Prototypes of label {\bf label} are added to all edges in the
  graph. The index of the label is returned. If the edges are
  already labeled with a label {\bf label}, only the index is
  returned.

\item{\bf int addGraphLabel(FFGraphLabel *label)}

Label {\bf label} is added to  the graph. The index of
the label is returned. If the graph is already labeled with a 
label {\bf label}, only the index is returned.

\item{\bf void delNodeLabel(int index, BOOL doFree=YES)}

The node labels with index {\bf index} are removed from all nodes. 
If {\bf doFree == YES} all labels are freed using the {\bf delete} 
operator.

\item{\bf void delEdgeLabel(int index, BOOL doFree=YES)}

The edge labels with index {\bf index} are removed from all edges. 
If {\bf doFree == YES} all labels are freed using the {\bf delete}
operator.

\item{\bf void delGraphLabel(int index, BOOL doFree=YES)}

The graph label with index {\bf index} is removed from the graph. 
If {\bf doFree == YES} the label is freed using the {\bf delete} 
operator.

\item{\bf void delNodeLabel(const String\& type, BOOL doFree=YES)}
  
  Removes the node label with type {\bf type} from all nodes in
  the graph. If {\bf doFree == YES} all labels are freed using
  the {\bf delete} operator.

\item{\bf void delEdgeLabel(const String\& type, BOOL doFree=YES)}
  
  Removes the edge label with type {\bf type} from all nodes in
  the graph. If {\bf doFree == YES} all labels are freed using
  the {\bf delete} operator.

\item{\bf void delGraphLabel(const String\& type, BOOL doFree=YES)}
  
  Removes the graph label with type {\bf type} from the graph.
  If {\bf doFree == YES} the label is freed using the {\bf
    delete} operator.
  
\item{\bf void assignGraphLabel(const FFGraphLabel *label, int
    i=-1)}
  
  Assigns {\bf label} to the graph label with index {\bf i} using
  the method {\bf void FFLabelClass::assign(FFLabelClass *)}.  If
  {\bf i == -1} the graph checks if it is labeled with a label of
  the same type as {\bf label} and assigns {\bf label} to it if
  it finds one.
  
\item{\bf void assignGraphLabel(const MPlex\verb|<|FFGraphLabel
    *\verb|>|\& set)}
  
  For all graph labels in {\bf set}, the graph checks if it is
  labeled with a label of the same type as the label from {\bf
    set} and assigns the label from {\bf set} to it if it finds
  one.
        
\item{\bf void setGraphLabel(FFGraphLabel *label , int index=-1,
    BOOL doFree=YES)}
  
  Sets the graph label of index {\bf index} to {\bf label}. If
  {\bf doFree == YES} the old label is freed using the {\bf
    delete} operator. If {\bf index == -1} the graph checks if
  it is labeled with a graph label of the same type as {\bf
    label} and sets it to {\bf label} if it finds one.
  
\item{\bf void setGraphLabel(const MPlex\verb|<|FFGraphLabel
    *\verb|>|\& set, BOOL doFree=YES)}
    
  For all graph labels in {\bf set}, the graph checks if it is
  labeled with a label of the same type as the label from {\bf
    set} and sets it to the label from {\bf set} if it finds one.
  If {\bf doFree == YES} the old label is freed using the {\bf
    delete} operator.
    
\item{\bf void setName(const String\& name)}

Sets the name of the graph to {\bf name}.

\item{\bf const String\& getName()}

Returns the name of the graph.
  
\item{\bf FFNodeClass *nodeWithValue(int k)}
  
  Returns the node with the value {\bf k}\footnote{compare
    FFNodeClass::setValue(int)}.
      
\item{\bf void rev()}

Reverts all edges.

\item{\bf void eliminateParallelEdges()}
  
  For all pairs of parallel edges, all but one are removed from
  the graph.  The edges to be removed are chosen randomly.

\item{\bf BOOL contains(FFNodeClass * n)}

Returns whether the graph contains node {\bf n}.

\item{\bf BOOL contains(FFEdgeClass * e)}

Returns whether the graph contains edge {\bf e}.

\item{\bf BOOL setStartNode(FFNode n)}
  
  Marks node {\bf n} as start-node. {\bf n} must be a node of the
  graph.  Unmarks the previous start-node.

\item{\bf FFNodeClass * getStartNode()}
  
  Returns the node that has been marked as start-node, or {\bf
    NULL} if no node has been marked as start-node.
  
\item{\bf BOOL setEndNodes(const AVLSet\verb|<|FFNodeClass
    *\verb|>|\& set)}
  
  Marks all nodes of {\bf set} as end-nodes. All nodes from {\bf
    set} must be nodes of the graph.

\item{\bf BOOL addEndNode(FFNode n)}
  
  Marks node {\bf n} as end-node. {\bf n} must be a node of the
  graph.

\item{\bf BOOL delEndNode(FFNode n)}
  
  Marks node {\bf n} as no longer being an end-node. {\bf n} must
  be a node of the graph.

\item{\bf void clearEndNodes()}

Marks all end-nodes as no longer being end-nodes.

\item{\bf const AVLSet\verb|<|FFNodeClass *\verb|>|\& getEndNodes()}
  
  Returns a set containing all nodes that are marked as being
  end-nodes.
  
\item{\bf int getNumNodes() const}

Returns the number of nodes in the graph.

\item{\bf int getNumEdges() const}    

Returns the number of edges in the graph.

\item{\bf FFGraph copy(FFGraph\& graph)}

Merges the graph into graph {\bf graph}. 

\item{\bf FFGraph copy()}

Returns a copy of the graph.

\item{\bf BOOL containsCycle()}

Returns whether the graph contains cycles.

\item{\bf BOOL weakConnected()}

Returns whether the graph is weak connected.

\item{\bf BOOL strongConnected()}

Returns whether the graph is strong connected.

\item{\bf BOOL isTree()}

Returns whether the graph is a tree.

\item{\bf void show()}

Prints the graph as ASCII to stdout.

\item{\bf void createGraphic(FFGraphWindow *w = NULL, FFGraphicNode 
	*gn=NULL, FFGraphicEdge *ge=NULL)}

Displays the graph in graphic context {\bf w}. If {\bf gn != NULL} it
is used instead of {\bf FFGraphicNode} to store the graphical
information of the nodes. If {\bf ge != NULL} it
is used instead of {\bf FFGraphicEdge} to store the graphical
information of the edges.

\item{\bf FFGraphWindow *getWindow() const}

Returns the graphic context the graph is displayed in, or {\bf NULL}
if the graph has no graphic context.

\item{\bf void removeGraphic()}

Removes the graph from the graphic context it is displayed in.
The graphic context is freed using the {\bf delete} operator.

\item{\bf FFGraphWindow *leaveWindow()}

Removes the graph from the graphic context it is displayed in.
The graphic context is returned.

\item{\bf FFGraphWindow *hideGraphic()}

Removes the graph from the graphic context it is displayed in.  All
graphic information in the nodes and edges remains unchanged. This
means the {\bf FFGraphicNode} and {\bf FFGraphicEdge} classes are not
deleted.  The graphic context is returned.

\end{itemize}

\subsubsection{FFNodeClass}

The class {\bf FFNodeClass} provides methods to manipulate nodes.
There is no public constructor for the class {\bf FFNodeClass}.
If you want a new node, use the method {\bf
  FFGraphClass::addNode(...)}.  {\bf FFNode} is a typedef to {\bf
  FFNodeClass *}

FFNodeClass defines the following methods:

\begin{itemize}

\item{\bf int getValue() const}

Returns the value of the node. Counterpart to {\bf setValue(...)}.

\item{\bf inline void setValue(int i)}
  
  Sets the value of the node to {\bf i}. Counterpart to {\bf
    getValue()}.  This method may be used by anyone and therefore
  the value is only very locally valid. {\bf This method must not
    be used by labels during init, parse or unparse!!!}

\item{\bf FFNode copy()}

Returns a copy of the node in a copy of the graph.
  
\item{\bf inline int indeg() const}

Returns the number of ingoing edges.

\item{\bf inline int outdeg() const}

Returns the number of outgoing edges.

\item{\bf void eliminateParallelEdges()}

For all parallel edges going into this node all but one are removed.

\item{\bf void assignLabel(FFNodeLabel *l, int index=-1)}

Assigns {\bf label} to the node label with index {\bf i} using
the method {\bf void FFLabelClass::assign(FFLabelClass *)}.
If {\bf i == -1} the node checks if it is labeled with a label
of the same type as {\bf label} and assigns {\bf label} to it
if it finds one.

\item{\bf void assignLabel(MPlex\verb|<|FFNodeLabel *\verb|>|\&
    set)}
  
  For all node labels in {\bf set} the node checks if it is
  labeled with a label of the same type as the label from {\bf
    set} and assigns the label from {\bf set} to it if it finds
  one.

\item{\bf void setLabel(FFNodeLabel *l, int i=-1, BOOL doFree=YES)}
  
  Sets the node label of index {\bf index} to {\bf label}. If
  {\bf doFree == YES} the old label is freed using the {\bf
    delete} operator. If {\bf index == -1} the node checks if it
  is labeled with a node label of the same type as {\bf label}
  and sets it to {\bf label} if it finds one.
  
\item{\bf void setLabel(MPlex\verb|<|FFNodeLabel *\verb|>|\& set,
    BOOL doFree=YES)}
    
  For all node labels in {\bf set}, the node checks if it is
  labeled with a label of the same type as the label from {\bf
    set} and sets it to the label from {\bf set} if it finds one.
  If {\bf doFree == YES} the old label is freed using the {\bf
    delete} operator.
 
\item{\bf int getLabelIndex(const String\& type,
    AVLSet\verb|<|int\verb|>| *set=NULL) const}
  
  If {\bf set == NULL} the index of the node label with type {\bf
    type} is returned, or -1 if none is found.  if {\bf set !=
    NULL} the indexes of all node labels which contain {\bf type}
  in their type are added to {\bf set}. The index of the first
  label that is found is returned, or -1, if none is found.

\item{\bf FFNodeLabel *getLabel(int index)}
  
  Returns the node label with index {\bf index}, or {\bf NULL} if
  the node has no label with index {\bf index}.
  
\item{\bf const MPlex\verb|<|FFNodeLabel *\verb|>| *getLabels()
    const}

Returns a MPlex containing all labels of the node.

\item{\bf FFGraph getGraph()}

Returns the graph the node is in.

\item{\bf void show()}

Prints the node to stdout.
  
\item{\bf AVLSet\verb|<|FFNode\verb|>| succ()}

Returns a set containing the successors of the node.

\item{\bf AVLSet\verb|<|FFNode\verb|>| pred()}

Returns a set containing the predecessor of the node.

\item{\bf const AVLSet\verb|<|FFEdge\verb|>|\& transToSucc() const}

Returns a set containing the edges to the successors of the node

\item{\bf const AVLSet\verb|<|FFEdge\verb|>|\& transFromPred()
    const}

Returns a set containing the edges from the predecessor of the node.

\item{\bf void succ(AVLSet\verb|<|FFNode\verb|>|\& set)}

Adds the successors of the node to {\bf set}.

\item{\bf void pred(AVLSet\verb|<|FFNode\verb|>|\& set)}

Adds the predecessor of the node to {\bf set}.
  
\item{\bf void setName(const String\& name)}

Sets the name of the node to {\bf name}.

\item{\bf const String\& getName() const}

Returns the name of the node.
  
\item{\bf void hideGraphic()}

Hides the node from a graphic context.

\item{\bf void showGraphic()}

Shows a node on the graphic context that has been hidden before.
  
\item{\bf FFGraphicNode *getGraphicNode() \\
    FFGraphicNode *operator()() \\
    FFGraphicNode *operator\verb|->|()}

Returns the {\bf FFGraphicNode} label, if the node is labeled with
{\bf FFGraphicNode}, or {\bf NULL} if it is not.

\end{itemize}

\subsubsection{FFEdgeClass}

The class {\bf FFEdgeClass} provides methods to manipulate edges.
There is no public constructor for the class {\bf FFEdgeClass}.
If you want a new Edge use the method {\bf
  FFGraphClass::addEdge(...)}. {\bf FFEdge} is a typedef to {\bf
  FFEdgeClass *}.  The class {\bf FFEdgeClass} provides the
following methods:

\begin{itemize}

\item{\bf FFEdge copy()}

Returns a copy of the edge in a copy of the graph.

\item{\bf void setValue(int k)}

Sets the value of the edge to k. Counterpart to {\bf getValue()}.
As anyone can use this method, values are only very locally valid.
{\bf Labels must not use this method during init, parse or unparse}.

\item{\bf int getValue()}

Returns the value of the edge. Counterpart to {\bf setValue(...)}.

\item{\bf void setWeight (FFCoord x)}

Sets the weight of the edge to {\bf x}.

\item{\bf FFCoord getWeight()}

Returns the weight of the edge.

\item{\bf void setName(const String\& name)}

Sets the name of the edge to {\bf name}.

\item{\bf const String\& getName()}

Returns the name of the edge.

\item{\bf FFGraph getGraph()}

Returns the graph the edge is in.

\item{\bf FFNode source()}

Returns the source node of the edge.

\item{\bf FFNode target()}

Returns the target node of the edge.

\item{\bf void revert()}

Reverts the edge.

\item{\bf void setTarget(FFNode n)}
  
  Sets the target of the edge to {\bf n}. {\bf n} must be a node
  of the same graph as the edge.

\item{\bf void setSource(FFNode n)}
  
  Sets the target of the edge to {\bf n}. {\bf n} must be a node
  of the same graph as the edge.

\item{\bf void assignLabel(FFEdgeLabel *label, int i=-1)}
  
  Assigns {\bf label} to the edge label with index {\bf i} using
  the method {\bf void FFLabelClass::assign(FFLabelClass *)}.  If
  {\bf i == -1} the edge checks if it is labeled with a label of
  the same type as {\bf label} and assigns {\bf label} to it if
  it finds one.

\item{\bf void setLabel(FFEdgeLabel *l, int i=-1, BOOL doFree=YES)}
  
  Sets the edge label of index {\bf index} to {\bf label}. If
  {\bf doFree == YES} the old label is freed using the {\bf
    delete} operator. If {\bf index == -1} the edge checks, if it
  is labeled with an edge label of the same type as {\bf label}
  and sets it to {\bf label} if it finds one.
  
\item{\bf void assignLabel(MPlex\verb|<|FFEdgeLabel *\verb|>|\&
    set)}
  
  For all edge labels in {\bf set} the edge checks if it is
  labeled with a label of the same type as the label from {\bf
    set} and assigns the label from {\bf set} to it if it finds
  one.
  
\item{\bf void setLabel(MPlex\verb|<|FFEdgeLabel *\verb|>|\& set,
    BOOL doFree=YES)}
  
  For all edge labels in {\bf set}, the edge checks if it is
  labeled with a label of the same type as the label from {\bf
    set} and sets it to the label from {\bf set} if it finds one.
  If {\bf doFree == YES} the old label is freed using the {\bf
    delete} operator.
 
\item{\bf int getLabelIndex(const String\& type,
    AVLSet\verb|<|int\verb|>| *set=NULL) const}
  
  If {\bf set == NULL} the index of the edge label with type {\bf
    type} is returned, or -1 if none is found.  if {\bf set !=
    NULL} the indexes of all edge labels which contain {\bf type}
  in their type are added to {\bf set}. The index of the first
  label that is found is returned, or -1, if none is found.

\item{\bf FFEdgeLabel *getLabel(int index)}
  
  Returns the edge label with index {\bf index}, or {\bf NULL} if
  the edge has no label with index {\bf index}.

\item{\bf MPlex\verb|<|FFEdgeLabel *\verb|>| *getLabels()}
  
  Returns a MPlex containing all labels of the edge.

\item{\bf void show()}
  
  Prints the edge to stdout.

\item{\bf void hideGraphic()}
  
  Hides the edge from a graphic context.

\item{\bf void showGraphic()}

Shows a previously hidden edge in the graphic context
  
\item{\bf FFGraphicEdge *getGraphicEdge() \\
    FFGraphicEdge *operator()() \\
    FFGraphicEdge *operator\verb|->|()}
  
  Returns the {\bf FFGraphicEdge} label if the edge is labeled
  with {\bf FFGraphicEdge} or {\bf NULL}, if it is not.

\end{itemize}
